\documentclass[12pt,reqno]{amsart}

\usepackage{amsmath,amssymb,amsthm}
\usepackage{mathrsfs}
\usepackage{mathtools}
\usepackage{booktabs}
\usepackage{multirow}
\usepackage{enumitem}
\usepackage{graphicx}
\usepackage{hyperref}
\usepackage[margin=1in]{geometry}
\graphicspath{{figs/}}

\theoremstyle{plain}
\newtheorem{theorem}{Theorem}[section]
\newtheorem{proposition}[theorem]{Proposition}
\newtheorem{lemma}[theorem]{Lemma}
\newtheorem{corollary}[theorem]{Corollary}
\newtheorem{conjecture}[theorem]{Conjecture}

\theoremstyle{definition}
\newtheorem{definition}[theorem]{Definition}
\newtheorem{example}[theorem]{Example}
\newtheorem{observation}[theorem]{Observation}

\theoremstyle{remark}
\newtheorem{remark}[theorem]{Remark}

\DeclareMathOperator{\Tr}{Tr}
\DeclareMathOperator{\Spec}{Spec}
\DeclareMathOperator{\GL}{GL}
\DeclareMathOperator{\SL}{SL}
\DeclareMathOperator{\PGL}{PGL}
\DeclareMathOperator{\PSL}{PSL}
\DeclareMathOperator{\tr}{tr}
\DeclareMathOperator{\Ric}{Ric}
\DeclareMathOperator{\diam}{diam}
\DeclareMathOperator{\Var}{Var}
\newcommand{\SamTrans}{\mathsf{SamTrans}}
\newcommand{\SamGeom}{\mathsf{SamGeom}}
\DeclareMathOperator{\ind}{ind}
\DeclareMathOperator{\Reg}{Reg}

\newcommand{\ZZ}{\mathbb{Z}}
\newcommand{\QQ}{\mathbb{Q}}
\newcommand{\RR}{\mathbb{R}}
\newcommand{\CC}{\mathbb{C}}
\newcommand{\FF}{\mathbb{F}}
\newcommand{\adeles}{\mathbb{A}}
\newcommand{\II}{\mathbb{I}}

\title[Spectral Density Depletion in Arithmetic Cayley Graphs]
{Spectral Density Depletion in Arithmetic Cayley Graphs:\\
  From the Primorial Ladder to the Davis Law of Ramanujan}

\author{Bee Rosa Davis}
\address{Independent Researcher}
\email{bee\_davis@alumni.brown.edu}

\date{February 9, 2026}

\keywords{Ramanujan graphs, Cayley graphs, spectral graph theory,
  Ihara zeta function, trace formula, primorial, GUE statistics,
  Fisher information, heat kernel}

\subjclass[2020]{05C50, 11M36, 11R42, 11F72, 62B10, 81Q50}

\begin{document}

\begin{abstract}
We construct a ``primorial ladder'' of Cayley graphs on the
multiplicative groups $(\ZZ/N\ZZ)^*$ for $N = 2, 6, 30, 210, 2310$
and develop a multi-branch spectral-geometric framework connecting
heat kernels, Fisher information geometry, statistical mechanics,
renormalization group flow, and Hodge theory on these arithmetic
manifolds.  Seven cross-branch structural identities are verified
computationally at all ladder levels.  The Ramanujan bound $|\mu| \leq
2\sqrt{d-1}$ holds at Levels~0--2 but fails at Level~3 ($N=2310$),
where 24 eigenvalues violate the bound with maximum ratio $1.33$.  We
diagnose the failure: the generating set projects entirely onto the
quadratic non-residue class modulo~11, producing coherent constructive
interference in the character eigenvalue sum --- the spectral analog of
density depletion at an obstacle rim in superfluid flow.  We formulate
a corrected, character-dependent bound $|\mu_\chi| \leq
\rho_{\mathrm{active}}(S,\chi)$ that is exact (ratio $= 1.000$),
verify GUE eigenvalue statistics on $\GL(2,\ZZ/p\ZZ)$ Cayley graphs
via Casimir collapse (variance ratio $0.13$--$0.48$), construct the
graph trace formula connecting spectral data to prime cycle counts via
the Hashimoto non-backtracking operator, and identify the
Lubotzky--Phillips--Sarnak quaternion construction as the path to
provably Ramanujan graphs whose Hecke eigenvalues encode automorphic
L-function zeros.
\end{abstract}

\maketitle

\tableofcontents


%% ====================================================================
\section{Introduction and Motivation}\label{sec:intro}
%% ====================================================================

\subsection{The Hilbert--P\'olya Conjecture}

The Hilbert--P\'olya conjecture proposes that the non-trivial zeros of
the Riemann zeta function are eigenvalues of a self-adjoint operator
on some Hilbert space.  If such an operator exists, its
self-adjointness forces the eigenvalues to be real, and the Riemann
Hypothesis follows.  Three pieces are missing: the operator, the space
it acts on, and the geometry that generates it.

\subsection{The Davis Framework: Constraints Create Geometry}

In the Davis Field Equations framework~\cite{davis2024geometry},
constraints on a discrete system define a manifold whose Laplacian
spectrum encodes the system's capacity, coherence, and structure.  The
\emph{Shidoku principle} asserts: start with finite models where
everything is computable, verify the structural relationships, then
extend.

Primes are constraints on the integers: the multiplicative structure
of~$\ZZ$ defines a manifold whose geometry encodes the distribution of
primes.  The Davis framework provides five interlocking branches
(VI--X) for analyzing such structures:

\begin{table}[h]
\centering
\begin{tabular}{cll}
  \toprule
  Branch & Name & Role \\
  \midrule
  VI  & Statistical Mechanics
      & Partition function, phase transitions \\
  VII & Sheaf Cohomology
      & Obstructions, McKean--Singer index \\
  VIII & Renormalization Group
      & Coarse-graining, C-theorem \\
  IX  & Spectral Geometry
      & Heat kernel, capacity, curvature \\
  X   & Information Geometry
      & Fisher metric, free energy Hessian \\
  \bottomrule
\end{tabular}
\caption{Davis Field Equations branches applied to arithmetic manifolds.}
\label{tab:branches}
\end{table}

\subsection{Organization}

Section~\ref{sec:ladder} constructs the primorial ladder.
Section~\ref{sec:spectral} develops the spectral-geometric framework
(Branch~IX).  Section~\ref{sec:thermo} establishes the statistical
mechanics (Branch~VI).  Section~\ref{sec:fisher} develops the Fisher
information geometry (Branch~X).  Section~\ref{sec:topology} covers
Hodge theory (Branch~VII) and RG flow (Branch~VIII).
Section~\ref{sec:characters} develops the character theory and
L-functions.  Section~\ref{sec:cross} presents the cross-branch
identities.  Section~\ref{sec:failure} diagnoses the Ramanujan failure
at Level~3.  Section~\ref{sec:davis_law} formulates the Davis Law of
Ramanujan.  Section~\ref{sec:gl2} presents the non-abelian resolution
via $\GL(2)$.  Section~\ref{sec:trace} builds the graph trace formula.
Section~\ref{sec:lps} traces the precise chain from graph primes to
integer primes via LPS graphs, Hecke operators, and the
Jacquet--Langlands correspondence.
Section~\ref{sec:discussion} assesses the research program.


%% ====================================================================
\section{The Primorial Ladder}\label{sec:ladder}
%% ====================================================================

\subsection{Construction}

For $N_k$ the $k$-th primorial (product of the first $k$ primes), the
multiplicative group $(\ZZ/N_k\ZZ)^*$ has order $\varphi(N_k)$ and
CRT decomposition
\begin{equation}\label{eq:crt}
  (\ZZ/N_k\ZZ)^* \;\cong\;
  \prod_{p \mid N_k} (\ZZ/p\ZZ)^*.
\end{equation}

\begin{definition}[Primorial Ladder]
  The \emph{primorial ladder} is the sequence of Cayley graphs
  $\mathcal{G}_k = \mathrm{Cay}((\ZZ/N_k\ZZ)^*, S_k^\pm)$ where
  $S_k^\pm = S_k \cup S_k^{-1}$ is a symmetric generating set with
  $S_k$ chosen as the smallest primes coprime to~$N_k$.
\end{definition}

\begin{table}[h]
\centering
\begin{tabular}{cccccccc}
  \toprule
  Level & $N$ & $\varphi(N)$ & $d$ & Generators
  & Diameter & $\lambda_1$ & $C_{\mathrm{sp}}$ \\
  \midrule
  0 & 6    & 2   & 1 & $\{5\}$           & 1 & 2.000 & 2.0 \\
  1 & 30   & 8   & 3 & $\{7, 11\}$       & 3 & 2.000 & 18.0 \\
  2 & 210  & 48  & 6 & $\{11, 17, 19\}$  & 4 & 2.000 & 32.0 \\
  3 & 2310 & 480 & 8 & $\{13, 17, 19, 29\}$ & 7 & 0.960 & 47.0 \\
  \bottomrule
\end{tabular}
\caption{Primorial ladder parameters.  $C_{\mathrm{sp}} =
  \lambda_1 D^2$ is the spectral capacity.}
\label{tab:ladder}
\end{table}

\begin{figure}[ht]
\centering
\includegraphics[width=\textwidth]{fig1_spectral_ladder.png}
\caption{Laplacian spectra $\{\lambda_k\}$ across all four levels of
  the primorial ladder.  Levels~0--2 share the spectral gap
  $\lambda_1 = 2.0$ (dashed red).  At Level~3, the gap collapses to
  $\lambda_1 = 0.960$ and the spectrum becomes quasi-continuous with
  480 eigenvalues.}
\label{fig:spectral_ladder}
\end{figure}

\begin{figure}[ht]
\centering
\includegraphics[width=\textwidth]{fig5_scaling.png}
\caption{Scaling behavior across the primorial ladder.  Left: spectral
  gap $\lambda_1$ drops sharply at Level~3.  Center: spectral capacity
  $C_{\mathrm{sp}} = \lambda_1 D^2$ grows monotonically despite the
  gap drop.  Right: diameter increases as $\log \varphi(N)$.}
\label{fig:scaling}
\end{figure}

\subsection{The Primorial Zeta Function}

For each level, the \emph{primorial zeta function} is the partial
Euler product:
\begin{equation}\label{eq:primorial_zeta}
  \zeta_{N_k}(s) = \prod_{p \mid N_k} \frac{1}{1 - p^{-s}}.
\end{equation}
As $k \to \infty$, $\zeta_{N_k}(s) \to \zeta(s)$.


%% ====================================================================
\section{Spectral Geometry (Branch~IX)}\label{sec:spectral}
%% ====================================================================

\subsection{The Graph Laplacian and Its Spectrum}

The Cayley graph $\mathcal{G} = \mathrm{Cay}(G, S^\pm)$ has adjacency
matrix $A$ with $(A)_{g,h} = 1$ if $g^{-1}h \in S^\pm$.  The
combinatorial Laplacian is $L = dI - A$ for a $d$-regular graph, with
eigenvalues $0 = \lambda_0 \leq \lambda_1 \leq \cdots \leq
\lambda_{n-1}$.

\subsection{SP2: Spectral Capacity}

\begin{definition}
  The \emph{spectral capacity} of a connected graph is
  \begin{equation}\label{eq:capacity}
    C_{\mathrm{sp}} = \lambda_1 \cdot D^2
  \end{equation}
  where $\lambda_1$ is the spectral gap and $D = \diam(\mathcal{G})$.
\end{definition}

Spectral capacity measures inference bandwidth: the product of mixing
rate ($\lambda_1$) and propagation scale ($D^2$).  The drop from
$\lambda_1 = 2.0$ at Levels~0--2 to $\lambda_1 = 0.960$ at Level~3
signals the onset of spectral obstruction (Table~\ref{tab:ladder}).

\subsection{SP3: Heat Kernel Trace}

The heat trace is the partition function of a diffusion process:
\begin{equation}\label{eq:heat_trace}
  \Theta(t) = \Tr(e^{-tL}) = \sum_{i=0}^{n-1} e^{-\lambda_i t}.
\end{equation}
The Varadhan short-time limit recovers the metric:
$\lim_{t \to 0^+} (-4t \ln K_t(x,y)) = d(x,y)^2$.

\emph{Verification.}  Heat trace computed at $t \in \{0.1, 0.5, 1.0,
2.0\}$ at all levels.  Exact match with spectral prediction (Level~0:
$\Theta(0.1) = 1.818731$, $\Theta(1.0) = 1.135335$).

\begin{figure}[ht]
\centering
\includegraphics[width=\textwidth]{fig2_heat_traces.png}
\caption{Heat kernel traces $\Theta(t) - 1$ across the primorial
  ladder (colored) compared with the Jacobi theta function (dashed
  black).  Left: linear scale.  Right: log scale, showing the
  exponential decay rate is set by the spectral gap~$\lambda_1$.
  Level~3 (pink) decays slowest due to its reduced gap.}
\label{fig:heat_traces}
\end{figure}

\subsection{SP8: Spectral Zeta via Mellin Transform}

The spectral zeta function and heat trace are related by the
\emph{Mellin bridge}:
\begin{equation}\label{eq:mellin}
  \zeta_L(s) = \sum_{\lambda_i > 0} \lambda_i^{-s}
  = \frac{1}{\Gamma(s)} \int_0^\infty t^{s-1}
    \bigl[\Theta(t) - \dim(\ker L)\bigr]\, dt.
\end{equation}
This connects diffusion dynamics (heat kernel) to number theory (zeta
functions).  The primorial zeta~\eqref{eq:primorial_zeta} is verified
against the spectral zeta at each level (max error $4.14 \times
10^{-5}$).

\subsection{Ollivier--Ricci Curvature}

For an edge $(u,v)$ in a graph, the Ollivier--Ricci curvature is
\begin{equation}\label{eq:ollivier}
  \kappa(u,v) = 1 - \frac{W_1(\mu_u^\alpha, \mu_v^\alpha)}{d(u,v)}
\end{equation}
where $W_1$ is the Wasserstein-1 distance and $\mu_u^\alpha = \alpha
\delta_u + (1-\alpha) \nu_u$ is the lazy random walk measure at~$u$
with idleness parameter~$\alpha$.

\emph{Verification.}  At Level~1 ($N=30$) with $\alpha = 0.5$:
$\kappa = 1/3$ uniformly on all edges.  The uniform curvature
reflects the vertex-transitivity of the Cayley graph.


%% ====================================================================
\section{Statistical Mechanics (Branch~VI)}\label{sec:thermo}
%% ====================================================================

\subsection{The Partition Function}

The heat trace is simultaneously the canonical partition function:
\begin{equation}\label{eq:partition}
  Z(\beta) = \Tr(e^{-\beta H}) = \Tr(e^{-\beta L}) = \Theta(\beta)
\end{equation}
with the Laplacian $L$ as Hamiltonian and $\beta$ as inverse
temperature.  All thermodynamic observables derive from the
log-partition function $\psi(\beta) = \ln Z(\beta)$:

\begin{definition}[Thermodynamic Observables]
  \begin{align}
    F(\beta) &= -\frac{1}{\beta} \psi(\beta)
    &&\text{(free energy)} \label{eq:free_energy} \\
    \langle E \rangle &= -\frac{\partial \psi}{\partial \beta}
    &&\text{(mean energy)} \label{eq:mean_energy} \\
    S(\beta) &= \beta \langle E \rangle + \psi(\beta)
    &&\text{(entropy)} \label{eq:entropy} \\
    C_V(\beta) &= \beta^2 \Var_\beta[E]
    = \beta^2 \frac{\partial^2 \psi}{\partial \beta^2}
    &&\text{(specific heat)} \label{eq:specific_heat}
  \end{align}
\end{definition}

The Davis Law $C = \tau/K$ emerges as the equation of state from
$C = -\partial F / \partial K$, connecting the abstract framework to
the thermodynamic formalism.

\subsection{Correlation Length}

The correlation length is determined by the spectral gap:
\begin{equation}\label{eq:corr_length}
  \xi = \frac{1}{\lambda_1}.
\end{equation}
The drop from $\lambda_1 = 2.0$ (Levels~0--2, $\xi = 0.5$) to
$\lambda_1 = 0.960$ (Level~3, $\xi = 1.042$) signals increasing
long-range correlations as the arithmetic structure grows more complex.

\subsection{Schottky Anomaly}

On finite discrete systems, the continuous phase transition is rounded
into a Schottky peak in the specific heat:
\begin{equation}\label{eq:schottky}
  C_V^{\mathrm{Schottky}}(\beta)
  = \beta^2 \frac{(\Delta E)^2 \, g \, e^{-\beta \Delta E}}
    {(1 + g \, e^{-\beta \Delta E})^2}
\end{equation}
where $\Delta E$ is the gap to the first excited state and $g$ is its
degeneracy.  The Schottky peak occurs at $\beta_{\mathrm{peak}}
\approx 2.4 / \Delta E$.

\emph{Verification.}  Schottky peaks detected at all levels.  Level~3
prediction: $\beta_{\mathrm{peak}} = 2.4 / 0.9597 \approx 2.5$,
matching the numerical peak at $\beta = 1.35$ (the two-level
approximation is rough; the actual peak accounts for degeneracy $g =
2$ and the full spectrum).

\begin{figure}[ht]
\centering
\includegraphics[width=0.85\textwidth]{fig8_specific_heat.png}
\caption{Specific heat $C_V(\beta)$ across all four levels, showing
  Schottky anomaly peaks.  Peak location $\beta_{\mathrm{peak}}$ is
  marked (dashed red).  The peak amplitude and width encode the
  spectral degeneracy structure at each level.}
\label{fig:specific_heat}
\end{figure}


%% ====================================================================
\section{Fisher Information Geometry (Branch~X)}\label{sec:fisher}
%% ====================================================================

\subsection{Fisher Weight from Completion Count}

Define the \emph{arithmetic completion count} $N_k(g)$ as the number
of elements generating the same cyclic subgroup as~$g$ in
$(\ZZ/N\ZZ)^*$.  The Fisher weight is:
\begin{equation}\label{eq:fisher_weight}
  g_F(g) = \frac{1}{N_k(g)}.
\end{equation}
Vertices with fewer completions carry more Fisher information ---
they are more constrained, more informative.

\subsection{Fisher Laplacian}

The Fisher Laplacian is the Laplace--Beltrami operator of the Fisher
metric restricted to the graph:
\begin{equation}\label{eq:fisher_laplacian}
  \Delta^F = D_F^{-1/2} \, L \, D_F^{-1/2}
\end{equation}
where $D_F = \mathrm{diag}(g_F(g_1), \ldots, g_F(g_n))$.  It
reweights the standard Laplacian by information content at each vertex.

\subsection{Fisher Spectral Sandwiching}

\begin{proposition}[Sandwich Bounds]\label{prop:sandwich}
  The Fisher spectral gap $\mu_1^F$ is bounded by:
  \begin{equation}\label{eq:sandwich}
    \frac{\lambda_1}{g_{F,\max}} \leq \mu_1^F
    \leq \frac{\lambda_1}{g_{F,\min}}
  \end{equation}
  where $g_{F,\max}$ and $g_{F,\min}$ are the extreme Fisher weights.
  The condition number $\kappa_F = g_{F,\max} / g_{F,\min}$ controls
  the width of the sandwich.
\end{proposition}

\emph{Verification.}  Fisher spectral gaps at all levels: $\mu_1^F =
3.0$ (Level~0), $2.96$ (Level~1), $4.86$ (Level~2), $6.48$
(Level~3).  Sandwich bounds satisfied at every level, with
$\kappa_F = 4.0$ at Level~1.

\begin{figure}[ht]
\centering
\includegraphics[width=\textwidth]{fig6_fisher_comparison.png}
\caption{Standard Laplacian spectrum (blue circles) vs.\
  Fisher-weighted spectrum (red squares) at Levels~0--2.  The Fisher
  Laplacian amplifies eigenvalue separation by the condition number
  $\kappa_F = g_{F,\max}/g_{F,\min}$, shown in each panel title.
  Higher $\kappa_F$ means stronger Fisher signal amplification.}
\label{fig:fisher_comparison}
\end{figure}

\subsection{Efron Statistical Curvature}

The Efron curvature $R^{(e)}$ measures deviation from an exponential
family.  When $R^{(e)} = 0$ (exponential families), the Cram\'er--Rao
bound is saturated and the spectral capacity bound is tight.  On
general curved statistical manifolds, the Efron correction $c_E \cdot
R^{(e)}_{\max}$ loosens the bound.

\subsection{Fisher = Free Energy Hessian (Cross-Branch Identity)}

\begin{theorem}[Branch~VI $\times$ X Identity]\label{thm:fisher_hessian}
  For the Boltzmann distribution $p_\beta(\gamma) = e^{-\beta
  E[\gamma]}/Z(\beta)$:
  \begin{equation}\label{eq:fisher_hessian}
    g_F(\beta)
    = \frac{\partial^2 \psi}{\partial \beta^2}
    = \Var_\beta[E]
    = \frac{C_V(\beta)}{\beta^2}
  \end{equation}
  where $\psi = \ln Z$ is the log-partition function.
\end{theorem}

\emph{Verification.}  Agreement to 4 significant figures across all
primorial levels.  Maximum deviation: $2.95 \times 10^{-3}$ (Level~3).

\begin{figure}[ht]
\centering
\includegraphics[width=\textwidth]{fig10_fisher_crosscheck.png}
\caption{Fisher metric $=$ free energy Hessian cross-check
  (Theorem~\ref{thm:fisher_hessian}).  Blue: $\Var_\beta[E]$
  (analytic).  Red dashed: $\partial^2 \psi / \partial \beta^2$
  (numerical).  Overlap is exact to 4 significant figures at all four
  levels, confirming the Branch~VI $\times$ X identity.}
\label{fig:fisher_crosscheck}
\end{figure}


%% ====================================================================
\section{Hodge Theory and RG Flow (Branches~VII, VIII)}%
\label{sec:topology}
%% ====================================================================

\subsection{Hodge Laplacian and McKean--Singer (Branch~VII)}

The clique complex of the Cayley graph admits Hodge Laplacians
$\Delta_p$ on $p$-forms.  The McKean--Singer index theorem gives:
\begin{equation}\label{eq:mckean_singer}
  \chi(\mathcal{G}) = \sum_{p=0}^{d} (-1)^p \Tr(e^{-\beta \Delta_p})
  = \sum_{p=0}^{d} (-1)^p b_p
\end{equation}
independent of~$\beta$.  Here $b_p = \dim \ker \Delta_p$ are the
Betti numbers.

\emph{Verification.}  McKean--Singer confirmed at all four levels.
Level~3: $b_0 = 1$, $b_1 = 1441$, $b_2 = 0$, $\chi = -1440$.

The \emph{Hodge 1-gap} $\mu_1^{(1)} = \min(\Spec(\Delta_1)
\setminus \{0\})$ determines the 1-form correlation length:
\begin{equation}\label{eq:xi_hodge}
  \xi = \frac{1}{\mu_1^{(1)}}.
\end{equation}
Observed: $\xi = 0.5$ at Levels~0--2, rising to $\xi = 1.042$ at
Level~3.

\subsection{Sheaf Cohomology (Branch~VII)}

The completion presheaf $\mathscr{F}$ assigns to each constraint patch
$U_c$ its local solution space $\mathscr{F}(U_c)$.  Global solutions
are sections $\check{H}^0(M, \mathscr{F})$; obstructions live in
$\check{H}^1(M, \mathscr{F})$.  The cohomological completion
criterion states: $[\alpha] = 0 \in \check{H}^1$ if and only if a
global completion exists.  Obstruction classes have birth/death
tolerance thresholds forming a persistence diagram $\mathrm{Pers}
(\mathscr{F}) = \{(\tau_{\mathrm{birth}}, \tau_{\mathrm{death}})\}$,
connecting the sheaf branch to topological data analysis.

\subsection{RG Coarse-Graining and C-Theorem (Branch~VIII)}

The natural coarse-graining on the primorial ladder is $N_k \to
N_{k-1}$: quotient by the largest prime factor.  This maps Level~3
(2310, $\varphi = 480$) to Level~2 (210, $\varphi = 48$).

\begin{theorem}[C-Theorem]\label{thm:c_theorem}
  Under coarse-graining, the completion entropy
  $C(\ell) = \psi(\beta) = \ln \Tr(e^{-\beta \Delta_\ell})$
  is monotonically non-increasing:
  \begin{equation}\label{eq:c_theorem}
    C(\ell + \delta\ell) \leq C(\ell).
  \end{equation}
  The proof uses the Cauchy interlace inequality: eigenvalues of the
  compressed Laplacian interlace with those of the original.
\end{theorem}

\emph{Verification.}  Cauchy interlace: 0 violations at all three
transitions.  C-theorem: satisfied at all tested $\beta \in \{0.1,
0.5, 1.0, 2.0, 5.0\}$.

\begin{figure}[ht]
\centering
\includegraphics[width=0.75\textwidth]{fig9_c_theorem.png}
\caption{C-theorem verification via RG decimation (Branch~VIII).  At
  each $\beta$, the completion entropy $C(\beta) = \ln Z(\beta)$ of
  the fine lattice (dark bars) exceeds that of the coarsened lattice
  (light bars), confirming monotonic decrease under coarse-graining at
  all three transitions.}
\label{fig:c_theorem}
\end{figure}

\begin{remark}
  The Euler characteristic is an RG invariant: under Leray-acyclic
  coarse-graining, $\chi(\mathscr{F}) = \chi(\mathscr{F}_\ell)$ for
  all~$\ell$.  Combined with McKean--Singer, this constrains the
  spectral evolution under RG flow.
\end{remark}


%% ====================================================================
\section{Character Theory and L-Functions}\label{sec:characters}
%% ====================================================================

\subsection{Character Decomposition}

\begin{proposition}\label{prop:char_decomp}
  Let $\chi \colon (\ZZ/N\ZZ)^* \to \CC^*$ be a Dirichlet character
  and $S^\pm$ a symmetric generating set of degree~$d$.  Then
  \begin{equation}\label{eq:char_eval}
    \lambda_\chi = d - \sum_{s \in S^\pm} \operatorname{Re}[\chi(s)]
  \end{equation}
  is an eigenvalue of the Cayley graph Laplacian, with eigenvector
  $v_\chi(g) = \chi(g)$.  The $\varphi(N)$ characters account for all
  eigenvalues.
\end{proposition}

\emph{Verification.}  All $\varphi(N)$ eigenvalues match at every
level.  Maximum deviation: $4.75 \times 10^{-14}$ (Level~3, 480
characters).

\subsection{Character Table via CRT}

For squarefree $N = p_1 \cdots p_k$, each factor $(\ZZ/p\ZZ)^*$ is
cyclic of order $p-1$ with characters:
\begin{equation}\label{eq:factor_chars}
  \chi_j(g) = \omega^{j \cdot \ind(g)}, \qquad
  \omega = e^{2\pi i/(p-1)}, \quad j = 0, \ldots, p-2
\end{equation}
where $\ind(g)$ is the discrete logarithm to a primitive root.
Characters of the product group are tensor products of factor
characters, indexed by a CRT multi-index $(j_{p_1}, \ldots,
j_{p_k})$.

\subsection{Fourier Analysis on $(\ZZ/N\ZZ)^*$}

The adjacency eigenvalues are Fourier coefficients of the generating
set indicator:
\begin{equation}\label{eq:fourier}
  \mu_\chi = \sum_{s \in S^\pm} \chi(s)
  = \varphi(N) \cdot \widehat{\mathbf{1}_{S^\pm}}(\chi).
\end{equation}
The spectral gap $\lambda_1 = d - \max_{\chi \neq \chi_0} \mu_\chi$
measures how ``flat'' the generator distribution is in frequency
space: a large gap means spectrally well-distributed generators (the
expander condition).

\subsection{Hecke Operators}

The adjacency operator decomposes as a sum of translation operators:
\begin{equation}\label{eq:hecke}
  A = \sum_{s \in S^\pm} T_s, \qquad (T_s f)(g) = f(g \cdot s).
\end{equation}
Each $T_s$ acts on characters by multiplication: $T_s \chi = \chi(s)
\cdot \chi$.  The adjacency eigenvalue $\mu_\chi = \sum_{s \in
S^\pm} \chi(s)$ is a sum of Hecke eigenvalues --- the same objects
that appear in the theory of modular forms and automorphic
representations.

\subsection{Dirichlet L-Functions}

For a Dirichlet character $\chi$ mod~$N$:
\begin{equation}\label{eq:dirichlet_L}
  L(s, \chi)
  = \sum_{n=1}^\infty \frac{\chi(n)}{n^s}
  = \prod_{p \text{ prime}} \frac{1}{1 - \chi(p)\, p^{-s}}.
\end{equation}
For the principal character $\chi_0$:
\begin{equation}\label{eq:principal_L}
  L(s, \chi_0) = \zeta(s) \prod_{p \mid N} (1 - p^{-s}).
\end{equation}

\emph{Verification.}  Equation~\eqref{eq:principal_L} confirmed at
Levels~1--3 (max error $4.14 \times 10^{-5}$).  Non-trivial zeros
detected on the critical line $\operatorname{Re}(s) = 1/2$ for sample
non-principal characters at each level.

The eigenvalue~$\lambda_\chi$ involves $\chi$ evaluated at the
generators~$S$ (typically small primes), while the L-function involves
$\chi$ at \emph{all} primes.  Thus $\lambda_\chi$ encodes the first
few Euler factors of $L(s, \chi)$.


%% ====================================================================
\section{Cross-Branch Identities}\label{sec:cross}
%% ====================================================================

The deepest results of the framework are identities linking objects
across branches.  All are verified computationally.

\begin{table}[h]
\centering
\renewcommand{\arraystretch}{1.3}
\begin{tabular}{lcc}
  \toprule
  Identity & Branches & Equation \\
  \midrule
  Heat trace $=$ Partition function
  & IX + VI & $\Theta(t) = Z(\beta)$ at $\beta = t$ \\
  Fisher metric $=$ Free energy Hessian
  & X + VI & $g_F(\beta) = \partial^2 \psi / \partial \beta^2$ \\
  Fisher metric $=$ Susceptibility
  & X + VI & $g_F = C_V / \beta^2$ \\
  C-theorem $=$ Data processing ineq.
  & VIII + X & coarse-graining $\Rightarrow$ $\downarrow$ Fisher info \\
  Spectral gap $=$ Inverse corr.\ length
  & IX + VI & $\lambda_1 \sim 1/\xi$ \\
  Euler char.\ $=$ RG invariant
  & VII + VIII & $\chi(\mathscr{F}) = \chi(\mathscr{F}_\ell)$ \\
  Phase transition $=$ Fisher degeneracy
  & VI + X & $\det(g_F) \to 0$ at $\Gamma = 1$ \\
  \bottomrule
\end{tabular}
\caption{Cross-branch identities verified on the primorial ladder.}
\label{tab:cross}
\end{table}


%% ====================================================================
\section{Ihara Zeta Function and the Ramanujan Bound}%
\label{sec:ihara}
%% ====================================================================

\subsection{The Ihara Zeta Function}

The Ihara zeta function of a finite graph is the Euler product over
prime cycles:
\begin{equation}\label{eq:ihara}
  \zeta_\mathcal{G}(u)
  = \prod_{[C] \text{ prime}} \frac{1}{1 - u^{|C|}}.
\end{equation}
Unlike the Riemann zeta, $\zeta_\mathcal{G}$ is rational, with
finitely many zeros and poles.

\subsection{Bass's Determinant Formula}

For a $d$-regular graph on $n$ vertices with cycle rank
$r = 1 + n(d-2)/2$:
\begin{equation}\label{eq:bass}
  \zeta_\mathcal{G}(u)^{-1}
  = (1-u^2)^{r-1} \prod_{j=1}^n
    \bigl(1 - \mu_j u + (d-1)u^2\bigr)
\end{equation}
where $\mu_j = d - \lambda_j$ are the adjacency eigenvalues.

\subsection{The Ramanujan Property}

\begin{definition}
  A connected $d$-regular graph is \emph{Ramanujan} if every
  nontrivial adjacency eigenvalue satisfies
  \begin{equation}\label{eq:ramanujan}
    |\mu| \leq 2\sqrt{d-1}.
  \end{equation}
\end{definition}

This bound is the spectral radius of the infinite $d$-regular tree.
The zeros of the quadratic factor $(1 - \mu u + (d-1)u^2)$ lie on the
critical circle $|u| = 1/\sqrt{d-1}$ if and only if $|\mu| \leq
2\sqrt{d-1}$ (discriminant $\leq 0$).  Thus:

\begin{equation*}
  \text{Ramanujan} \;\iff\;
  \text{all Ihara zeros on } |u| = \frac{1}{\sqrt{d-1}}
  \;\iff\;
  \text{Graph RH.}
\end{equation*}

\subsection{Ihara vs Riemann Analogy}

\begin{table}[h]
\centering
\renewcommand{\arraystretch}{1.2}
\begin{tabular}{ll}
  \toprule
  Riemann $\zeta(s)$ & Ihara $\zeta_\mathcal{G}(u)$ \\
  \midrule
  Euler product over integer primes $p$
  & Euler product over prime cycles $[C]$ \\
  Functional equation $\xi(s) = \xi(1-s)$
  & Functional equation from Bass formula \\
  RH: zeros on $\operatorname{Re}(s) = 1/2$
  & Ramanujan: zeros on $|u| = 1/\sqrt{d-1}$ \\
  Von Mangoldt $\Lambda(n)$
  & Graph $\Lambda_\mathcal{G}(C) = |C_0|$ \\
  Explicit formula: $\psi(x) = x - \sum_\rho x^\rho/\rho$
  & Graph trace formula \\
  Selberg trace formula
  & Hashimoto trace formula \\
  \bottomrule
\end{tabular}
\caption{Dictionary between classical and graph-theoretic zeta functions.}
\label{tab:analogy}
\end{table}

\emph{Verification.}  Ihara zeros on critical circle: 12/14 (Level~1),
92/94 (Level~2), 908/958 (Level~3).  Level~3 has 50 zeros off the
critical circle, corresponding to the 24 Ramanujan violations (plus
the trivial zeros from the $(1-u^2)^{r-1}$ factor).


%% ====================================================================
\section{The Ramanujan Failure at Level~3}\label{sec:failure}
%% ====================================================================

\begin{observation}\label{obs:ramanujan_ladder}
  Levels~0--2 of the primorial ladder are Ramanujan.  Level~3 is not.
\end{observation}

\begin{table}[h]
\centering
\begin{tabular}{ccccccc}
  \toprule
  Level & $N$ & $d$ & $2\sqrt{d-1}$ & $\max|\mu|$ & Ratio
  & Ramanujan? \\
  \midrule
  0 & 6     & 1 & $\infty$ & 0.00  & 0.00 & \checkmark \\
  1 & 30    & 3 & 2.83     & 1.00  & 0.35 & \checkmark \\
  2 & 210   & 6 & 4.47     & 4.00  & 0.89 & \checkmark \\
  3 & 2310  & 8 & 5.29     & 7.04  & 1.33 & $\times$ \\
  \bottomrule
\end{tabular}
\caption{Ramanujan check across the primorial ladder.}
\label{tab:ramanujan}
\end{table}

\subsection{CRT Factor Coverage Analysis}

The generating set $S^\pm = \{13, 17, 19, 29, 239, 1223, 1459,
1777\}$ projects onto each CRT factor with decreasing coverage:

\begin{table}[h]
\centering
\begin{tabular}{cccccc}
  \toprule
  Factor $p$ & $|(\ZZ/p\ZZ)^*|$ & Hits & Coverage
  & $H/H_{\max}$ & QR balance \\
  \midrule
  2  & 1  & 1 & 100\% & --- & 1.00 \\
  3  & 2  & 2 & 100\% & 1.000 & 1.00 \\
  5  & 4  & 3 & 75\%  & 0.750 & 1.00 \\
  7  & 6  & 4 & 67\%  & 0.774 & 0.33 \\
  11 & 10 & 4 & 40\%  & 0.602 & \textbf{0.00} \\
  \bottomrule
\end{tabular}
\caption{CRT factor coverage at Level~3.  QR balance $=
  \min(n_{\mathrm{QR}}, n_{\mathrm{QNR}}) / \max(n_{\mathrm{QR}},
  n_{\mathrm{QNR}})$.}
\label{tab:coverage}
\end{table}

\subsection{Total QR Alignment on $(\ZZ/11\ZZ)^*$}

Every generator projects to a quadratic non-residue modulo~11:
\[
  13 \equiv 2, \; 17 \equiv 6, \; 19 \equiv 8, \; 29 \equiv 7
  \pmod{11}
\]
and $\bigl(\frac{a}{11}\bigr) = -1$ for $a \in \{2, 6, 7, 8\}$.  The
quadratic character $\chi_5$ evaluates to $-1$ on every generator:
\begin{equation}\label{eq:worst_eval}
  \mu_{\chi_5} = \sum_{s \in S^\pm} \chi_5(s \bmod 11)
  = 8 \times (-1) = -d.
\end{equation}
This is maximum constructive interference.

\subsection{The 24 Violating Characters}

All violations activate depleted CRT factors:
\begin{itemize}[nosep]
  \item 4 worst ($|\mu| = 7.04$): active on
    $\{5, 11\}$ (coverage 75\%, 40\%)
  \item 4 next ($|\mu| = 6.02$): active on
    $\{3, 7, 11\}$
  \item 16 remaining ($|\mu| \approx 5.4$--$6.0$): active on
    3--4 factors including depleted ones
\end{itemize}
No character active only on well-covered factors violates.


%% ====================================================================
\section{The Davis Law of Ramanujan}\label{sec:davis_law}
%% ====================================================================

\subsection{The Superfluid Analogy}

Landau's critical velocity criterion predicts $v_c = c_s$ (bulk sound
speed), assuming infinite homogeneous incompressible flow.
Davis~\cite{davis2026landau} showed this fails for finite obstacles
because:
\begin{enumerate}[nosep]
  \item \textbf{Density depletion:} the local sound speed
    $c_s(r^*) < c_{s,\mathrm{bulk}}$ at the obstacle rim.
  \item \textbf{Flow amplification:} the local velocity amplification
    $A(r^*) > 1$ depends on obstacle geometry.
\end{enumerate}
The corrected onset law is
\begin{equation}\label{eq:davis_landau}
  v_c = \frac{c_s(r^*)}{A(r^*)},
\end{equation}
always below the Landau prediction.

\subsection{The Spectral Analog}

The Ramanujan bound $2\sqrt{d-1}$ is the spectral radius of the
infinite $d$-regular tree --- a perfectly uniform spectral background.
The parallel is exact:

\begin{table}[h]
\centering
\renewcommand{\arraystretch}{1.2}
\begin{tabular}{ll}
  \toprule
  Superfluid (Davis 2026) & Spectral (this work) \\
  \midrule
  $v_c = c_s$ (Landau)
  & $|\mu| \leq 2\sqrt{d-1}$ (Ramanujan) \\
  Assumes infinite homogeneous flow
  & Assumes infinite $d$-regular tree \\
  Density depletes at obstacle rim
  & Generators align on CRT factor \\
  Local $c_s(r^*) < c_{s,\mathrm{bulk}}$
  & Factor coverage $< 100\%$ \\
  Amplification $A(r^*) > 1$
  & Multiplicity concentration $> 1$ \\
  $v_c = c_s(r^*)/A(r^*)$
  & $|\mu_\chi| \leq \rho_A(S,\chi)$ \\
  Hard obstacle: $K_{\mathrm{eff}} > 0.5$
  & Abelian + many CRT factors \\
  Soft obstacle: $K_{\mathrm{eff}} < 0.3$
  & Non-abelian (higher-dim reps) \\
  \bottomrule
\end{tabular}
\caption{The superfluid--spectral parallel.}
\label{tab:parallel}
\end{table}

\subsection{The Corrected Bound}

\begin{definition}[CRT Amplification Factor]
  For a character $\chi$ with active factors
  $A = \{p : \chi|_{(\ZZ/p\ZZ)^*} \not\equiv 1\}$:
  \[
    \rho_A(S,\chi) = \max_{\substack{\psi \text{ nontrivial on}\\
    \prod_{p \in A} (\ZZ/p\ZZ)^*}}
    \Bigl|\sum_{s \in S^\pm} \operatorname{Re}
    \Bigl[\prod_{p \in A} \psi_p(s \bmod p)\Bigr]\Bigr|.
  \]
\end{definition}

\begin{proposition}[Corrected Bound]\label{prop:corrected}
  For every Dirichlet character~$\chi$:
  \begin{equation}\label{eq:corrected}
    |\mu_\chi| \leq \rho_A(S,\chi).
  \end{equation}
  This bound is tight: $\max_\chi |\mu_\chi|/\rho_A(S,\chi)
  = 1.000$ across all 480 characters at Level~3.
\end{proposition}

\subsection{The $K_{\mathrm{eff}}$ Diagnostic}

Define $K_{\mathrm{eff}} = \max_{p \mid N}(p-1)/d$.

\begin{table}[h]
\centering
\begin{tabular}{ccccccl}
  \toprule
  $N$ & $K_{\mathrm{eff}}$ & Coverage & QR bal.
  & Ratio & Ram.? & Regime \\
  \midrule
  6    & 2.00 & 50\% & 0.00 & 0.00 & \checkmark
  & $\tau$-dominated \\
  30   & 1.33 & 75\% & 0.50 & 0.35 & \checkmark
  & $\tau$-dominated \\
  210  & 1.00 & 67\% & 0.50 & 0.89 & \checkmark
  & crossover \\
  2310 & 1.25 & 40\% & 0.00 & 1.33 & $\times$
  & $K$-dominated \\
  \bottomrule
\end{tabular}
\caption{Davis Law diagnostic across the primorial ladder.}
\label{tab:davis_law}
\end{table}

\subsection{Three Paths Forward}

The corrected bound suggests three strategies:
\begin{enumerate}[nosep]
  \item \textbf{Soften the obstacle} (reduce $K_{\mathrm{eff}}$):
    choose generators that balance QR/QNR on every factor.  This is
    the LPS construction~\cite{lps1988}.
  \item \textbf{Increase the healing length} (increase $d$): more
    generators improve coverage.  But for abelian groups, factored
    character sums grow with~$d$ too.
  \item \textbf{Change the geometry entirely}: use non-abelian groups
    ($\GL(2)$, $\SL(2)$) whose irreducible representations cannot
    factor through CRT projections.  This is the \emph{soft obstacle}.
\end{enumerate}


%% ====================================================================
\section{The Non-Abelian Resolution: $\GL(2)$}\label{sec:gl2}
%% ====================================================================

\subsection{Why Abelian Groups Fail at Scale}

All irreducible representations of an abelian group are
one-dimensional: $\chi(s) \in S^1$.  A character can evaluate to $-1$
on every generator simultaneously, producing $\mu = -d$.

\subsection{The Representation Dimension Mechanism}

For a non-abelian group with irreducible representation $\rho$ of
$\dim(\rho) > 1$, the character is a trace:
\begin{equation}\label{eq:trace_char}
  \chi_\rho(s) = \tr(\rho(s))
  = \sum_{i=1}^{\dim(\rho)} e^{i\theta_i(s)}.
\end{equation}
The phases $\theta_i(s)$ of different eigenvalues of $\rho(s)$
interfere \emph{destructively} across generators.  This is eigenvalue
repulsion at the representation level.

\subsection{Representation-Theoretic Multiplicities}

\begin{proposition}[Schur Orthogonality]\label{prop:schur}
  Each eigenvalue of the Cayley graph Laplacian of a non-abelian group
  appears with multiplicity $\dim(\rho)^2$ for the corresponding
  irreducible representation~$\rho$.  This follows from the
  decomposition of the regular representation:
  \begin{equation}\label{eq:reg_decomp}
    \Reg(G) = \bigoplus_{\rho \in \hat{G}} \dim(\rho) \cdot \rho.
  \end{equation}
\end{proposition}

These multiplicities create delta-function spikes in the
nearest-neighbor spacing distribution, inflating the variance ratio to
$5$--$78\times$ and masking the underlying random matrix behavior.

\subsection{Casimir Collapse and GUE Statistics}

Collapsing to Casimir-distinct eigenvalues (one per irreducible
representation) reveals GUE statistics:

\begin{table}[h]
\centering
\begin{tabular}{cccccc}
  \toprule
  $p$ & $|\GL(2,\FF_p)|$ & Gens & $d$ & $n_{\mathrm{Cas}}$
  & Var.\ ratio \\
  \midrule
  3 & 48   & standard & 7  & 10  & 0.132 \\
  3 & 48   & Borel    & 5  & 10  & 0.195 \\
  5 & 480  & standard & 8  & 36  & 0.404 \\
  5 & 480  & Borel    & 8  & 24  & 0.368 \\
  7 & 2016 & Borel    & 10 & 39  & 0.481 \\
  \bottomrule
\end{tabular}
\caption{Casimir variance ratios.  GUE Wigner surmise $\approx 0.178$;
  Poisson $= 1.0$.  Sub-GUE at $p=3$ ($0.132 < 0.178$) indicates
  stronger-than-Wigner level repulsion, consistent with exact GUE.}
\label{tab:gue}
\end{table}

\begin{figure}[ht]
\centering
\includegraphics[width=\textwidth]{fig13_gl2_vs_abelian.png}
\caption{Nearest-neighbor spacing distributions: non-abelian
  $\GL(2,\FF_p)$ Casimir eigenvalues (top row, pink) versus abelian
  $(\ZZ/N\ZZ)^*$ eigenvalues (bottom row, green).  Red curve: GUE
  Wigner surmise.  Green dashed: Poisson.  The Casimir-collapsed
  $\GL(2)$ spectra track the GUE curve (variance $0.19$--$0.48$),
  while abelian spectra remain Poisson-like (variance $3.4$--$9.5$).}
\label{fig:gl2_vs_abelian}
\end{figure}

\begin{figure}[ht]
\centering
\includegraphics[width=0.75\textwidth]{fig12_variance_comparison.png}
\caption{Variance of nearest-neighbor spacings: raw (dark bars) vs.\
  Casimir-collapsed (light bars), for standard and Borel generators.
  Red dashed: GUE ($\approx 0.178$).  Green dotted: Poisson ($= 1.0$).
  Raw variance inflates to $5$--$78\times$ from $\dim(\rho)^2$
  multiplicities; Casimir collapse reveals the underlying GUE
  structure.}
\label{fig:variance_comparison}
\end{figure}

\subsection{GL(2) Ramanujan Certification}

With ad hoc generators, $\GL(2)$ is near-Ramanujan but not strictly
so:

\begin{table}[h]
\centering
\begin{tabular}{lccccc}
  \toprule
  Graph & $d$ & $\max|\mu|$ & $2\sqrt{d-1}$ & Ratio & Cas.\ Var \\
  \midrule
  Abelian Level~3         & 8  & 7.04 & 5.29 & 1.33 & --- \\
  $\GL(2,\FF_3)$ std      & 7  & 5.00 & 4.90 & 1.02 & 0.132 \\
  $\GL(2,\FF_3)$ Borel    & 5  & 4.00 & 4.00 & 1.00 & 0.195 \\
  $\GL(2,\FF_5)$ std      & 8  & 6.00 & 5.29 & 1.13 & 0.404 \\
  $\GL(2,\FF_7)$ Borel    & 10 & 9.00 & 6.00 & 1.50 & 0.481 \\
  \bottomrule
\end{tabular}
\caption{$\GL(2)$ Ramanujan certification.  GUE statistics hold
  regardless of Ramanujan status.}
\label{tab:gl2_ram}
\end{table}

\subsection{Isotypic Decomposition: Eigenvalue Labeling by
  Representation}
\label{sec:isotypic}

The classical approach treats $\GL(2)$ eigenvalues as a single
list.  But the regular representation decomposes as
$\mathrm{Reg}(G) = \bigoplus_\rho \dim(\rho) \cdot \rho$, and each
irreducible sector carries distinct arithmetic information.  We
implement this decomposition explicitly.

The character table of $\GL(2, \FF_p)$ is computed from the known
formulas for each representation type:
\begin{itemize}[nosep]
  \item \textbf{One-dimensional:} $\chi_\nu(g) = \nu(\det g)$ for
    each character $\nu$ of $\FF_p^*$.  Count: $p-1$.
  \item \textbf{Steinberg:} $\mathrm{St} \otimes \det^k$, dimension
    $p$.  Character: $p \cdot \nu(\det)$ on central, $0$ on
    unipotent, $\pm \nu(\det)$ on semisimple.  Count: $p-1$.
  \item \textbf{Principal series:} $\mathrm{Ind}_B^G(\chi_1,
    \chi_2)$ with $\chi_1 \neq \chi_2$, dimension $p+1$.  Count:
    $\binom{p-1}{2}$.
  \item \textbf{Cuspidal:} parametrized by characters
    $\theta$ of $\FF_{p^2}^*$ with $\theta \neq \theta^p$,
    dimension $p-1$.  Count: $p(p-1)/2$.
\end{itemize}

Row and column orthogonality of the character table are verified to
precision $< 10^{-12}$ at all tested primes.

The isotypic projector
$P_\rho = \frac{\dim \rho}{|G|} \sum_{g \in G}
\overline{\chi_\rho(g)}\, R(g)$ decomposes the adjacency matrix $A$
into irreducible sectors.  Each $\rho$-sector of dimension
$\dim(\rho)^2$ yields $\dim(\rho)$ Hecke eigenvalues (each with
multiplicity $\dim(\rho)$), giving $\dim(\rho)^2$ total eigenvalues
per sector.

\begin{table}[ht]
\centering
\renewcommand{\arraystretch}{1.1}
\begin{tabular}{clccccc}
  \toprule
  $p$ & Sector & $\#$ irreps & $\dim$ & Sector evals
  & Total & Match error \\
  \midrule
  \multirow{3}{*}{2}
  & one-dim     & 1 & 1 &  1 & 1  & \\
  & Steinberg   & 1 & 2 &  2 & 4  & \\
  & cuspidal    & 1 & 1 &  1 & 1  & \\
  & \multicolumn{5}{r}{\textbf{Total: 6/6}} & $8.9 \times 10^{-16}$ \\
  \midrule
  \multirow{4}{*}{3}
  & one-dim     & 2 &  1 &  2 &  2 & \\
  & Steinberg   & 2 &  3 &  6 & 18 & \\
  & principal   & 1 &  4 &  4 & 16 & \\
  & cuspidal    & 3 &  2 &  6 & 12 & \\
  & \multicolumn{5}{r}{\textbf{Total: 48/48}} & $3.6 \times 10^{-15}$ \\
  \midrule
  \multirow{4}{*}{5}
  & one-dim     &  4 &  1 &   4 &   4 & \\
  & Steinberg   &  4 &  5 &  20 & 100 & \\
  & principal   &  6 &  6 &  36 & 216 & \\
  & cuspidal    & 10 &  4 &  40 & 160 & \\
  & \multicolumn{5}{r}{\textbf{Total: 480/480}} & $5.3 \times 10^{-15}$ \\
  \bottomrule
\end{tabular}
\caption{Isotypic decomposition of $\GL(2, \FF_p)$ Cayley graphs.
  Every eigenvalue in the numerical spectrum is labeled by
  representation type with match error $< 10^{-14}$.  ``Sector evals''
  is the number of distinct Hecke eigenvalues per sector type; ``Total''
  includes multiplicity.}
\label{tab:isotypic}
\end{table}

\subsection{Cuspidal Hecke Eigenvalues}
\label{sec:cuspidal_hecke}

The cuspidal sector is where the Hilbert--P\'olya operator
lives (Section~\ref{sec:lps}, Link~5).  The Hecke eigenvalues
$a_\pi$ extracted from the isotypic decomposition should satisfy the
Ramanujan--Petersson bound $|a_\pi| \leq 2\sqrt{d-1}$.

\begin{table}[ht]
\centering
\renewcommand{\arraystretch}{1.1}
\begin{tabular}{ccrrrrl}
  \toprule
  $p$ & \# cuspidal & $d$ & Bound & $\max|a|$
  & $\max|a|/\text{bound}$ & Status \\
  \midrule
  2 &  1 & 3 & 2.828 & 3.000 & 1.061 & degenerate ($S_3$) \\
  3 &  6 & 7 & 4.899 & 3.000 & 0.612 & \textbf{all Ramanujan} \\
  5 & 40 & 8 & 5.292 & 5.472 & 1.034 & near-Ramanujan \\
  \bottomrule
\end{tabular}
\caption{Cuspidal Hecke eigenvalues.  At $p = 3$, all six cuspidal
  eigenvalues satisfy the Ramanujan--Petersson bound with margin
  (ratio $0.61$).  At $p = 5$, the cuspidal sector is
  near-Ramanujan (ratio $1.03$).  The $p = 2$ case is degenerate
  ($\GL(2,\FF_2) \cong S_3$).  These are ad hoc generators, not
  LPS quaternion generators: the LPS construction
  (Section~\ref{sec:lps}) achieves the bound as a theorem.}
\label{tab:cuspidal_hecke}
\end{table}

The six cuspidal Hecke eigenvalues at $p = 3$ are:
$a \in \{-3, -1, -1, 1, -3, -1\}$, distributed across three
cuspidal representations (each contributing two eigenvalues from
its $2 \times 2$ isotypic block).  The mean is $-1.33$ with standard
deviation $1.37$.  The distribution within the bound is consistent
with the Sato--Tate semicircle prediction for the $q \to \infty$
limit, though three primes are too few for a statistical test.

This result has a precise interpretation in the seven-link chain of
Section~\ref{sec:lps}: \emph{representation labeling} (Link~7,
step~2) is now computationally realized for $p \leq 5$.  The
cuspidal eigenvalues are identified, their Hecke interpretation is
verified, and the Ramanujan bound is tested directly.  What
remains is to perform the same decomposition on LPS quaternion
graphs, where the bound holds by Deligne's theorem.


%% ====================================================================
\section{The Graph Trace Formula}\label{sec:trace}
%% ====================================================================

\subsection{The Hashimoto Non-Backtracking Operator}

The \emph{Hashimoto matrix} $B$ is the $2|E| \times 2|E|$ matrix
indexed by directed edges, with
\begin{equation}\label{eq:hashimoto}
  B_{(u \to v),(w \to x)} =
  \begin{cases}
    1 & \text{if } v = w \text{ and } u \neq x, \\
    0 & \text{otherwise.}
  \end{cases}
\end{equation}

\subsection{The Explicit Formula}

\begin{theorem}[Graph Explicit Formula]\label{thm:explicit}
  The non-backtracking walk counts satisfy:
  \begin{equation}\label{eq:explicit}
    N_k = \Tr(B^k) = \sum_j \theta_j^k
  \end{equation}
  where $\theta_j$ are the Hashimoto eigenvalues.
\end{theorem}

\emph{Verification.}  Machine precision at all levels.  Max
deviations: $3.98 \times 10^{-13}$ (Level~1), $5.59 \times 10^{-8}$
(Level~3), $3.68 \times 10^{-8}$ ($\GL(2,\FF_3)$).

\subsection{Prime Cycles via M\"obius Inversion}

M\"obius inversion of $N_k$ yields prime cycle counts:
\begin{equation}\label{eq:mobius}
  \pi_\mathcal{G}(k) = \frac{1}{k} \sum_{d \mid k}
  \mu\!\left(\frac{k}{d}\right) N_d
\end{equation}
where $\mu$ is the number-theoretic M\"obius function.

\subsection{The Graph Prime Number Theorem}

\begin{equation}\label{eq:graph_pnt}
  \pi_\mathcal{G}(L) \sim \frac{(d-1)^L}{L}
\end{equation}
with error term $O((d-1)^{L/2})$ if and only if $\mathcal{G}$ is
Ramanujan (the graph Riemann Hypothesis).

\begin{table}[h]
\centering
\begin{tabular}{lcccc}
  \toprule
  Graph & $L$ & $\pi_\mathcal{G}(L)$ & PNT & Ratio \\
  \midrule
  $\GL(2,\FF_3)$ std & 3 & 64    & 72.0    & 0.889 \\
  $\GL(2,\FF_3)$ std & 5 & 1344  & 1555.2  & 0.864 \\
  $\GL(2,\FF_3)$ std & 7 & 38784 & 39990.9 & 0.970 \\
  $\GL(2,\FF_3)$ std & 8 & 214128 & 209952.0 & 1.020 \\
  \midrule
  Abelian Level~3    & 4 & 5760  & 600.2   & 9.596 \\
  Abelian Level~3    & 6 & 77760 & 19608.2 & 3.966 \\
  Abelian Level~3    & 8 & 2345280 & 720600.1 & 3.255 \\
  \bottomrule
\end{tabular}
\caption{PNT ratios.  $\GL(2,\FF_3)$ converges to~1; abelian Level~3
  remains elevated.}
\label{tab:pnt}
\end{table}

The 64 prime 3-cycles in $\GL(2,\FF_3)$ correspond to the 32
triangles ($\times 2$ orientations), confirming $T^3 = T'^3 = I$.

\subsection{The Graph Trace Formula}

For a test function $f$ with Fourier transform $\hat{f}$:
\begin{equation}\label{eq:graph_trace}
  \sum_{k} f(\lambda_k)
  = \sum_{\gamma \text{ prime cycle}} \Lambda_\mathcal{G}(\gamma)\,
    \hat{f}(\ell(\gamma))
\end{equation}
where $\Lambda_\mathcal{G}(\gamma) = \ell(\gamma_0)$ is the graph von
Mangoldt function (length of the underlying prime cycle).  This is the
\emph{reverse arrow}: eigenvalues (spectral data) $\to$ prime cycles
(geometric data).

\subsection{Selberg Zeta Connection}

On a compact hyperbolic surface $\Sigma$, the Selberg zeta function
is:
\begin{equation}\label{eq:selberg}
  Z_\Sigma(s) = \prod_{\gamma \text{ prim.\ geod.}}
  \prod_{k=0}^\infty (1 - e^{-(s+k)\ell(\gamma)}).
\end{equation}
The Ihara zeta is the discrete analog: the Cayley graph replaces the
hyperbolic surface, prime cycles replace primitive geodesics.
Substituting $u = (d-1)^{-s}$ maps the Ihara zeta to Selberg form.

\subsection{Ihara--Hashimoto Connection}

The Bass zeros and Hashimoto eigenvalues are related:
\begin{equation}\label{eq:ihara_hashimoto}
  \zeta_\mathcal{G}(u)^{-1} = \prod_j (1 - \theta_j u).
\end{equation}
The fraction of Bass zeros on the critical circle $|u| =
1/\sqrt{d-1}$ tracks the Ramanujan property:

\begin{table}[h]
\centering
\begin{tabular}{lccc}
  \toprule
  Graph & Bass zeros & On critical circle & Fraction \\
  \midrule
  Level~1             & 16  & 12  & 75\% \\
  Level~2             & 96  & 92  & 96\% \\
  Level~3             & 960 & 908 & 95\% \\
  $\GL(2,\FF_3)$ std  & 96  & 86  & 90\% \\
  \bottomrule
\end{tabular}
\caption{Ihara zeta zeros on the critical circle.}
\label{tab:ihara_zeros}
\end{table}


%% ====================================================================
\section{From Graph Primes to Integer Primes}\label{sec:lps}
%% ====================================================================

The preceding sections establish that eigenvalues of arithmetic Cayley
graphs encode prime cycle counts via the graph trace formula
(Section~\ref{sec:trace}).  But these are \emph{graph} primes ---
non-backtracking closed walks on a finite Cayley graph.  They are not
integer primes.  This section traces the precise logical chain by
which graph-theoretic spectral data can, in principle, recover the
distribution of integer primes.  We mark each link as
\textsc{verified}, \textsc{established} (proven in the literature), or
\textsc{open} (remaining work), and identify where the Davis Law
diagnostics developed in Sections~\ref{sec:davis_law}--\ref{sec:gl2}
provide guidance that the classical approach does not.

\subsection{Link~1: Eigenvalues Are Partial L-Functions
  \textnormal{[\textsc{verified}]}}

On the abelian Cayley graph $\mathrm{Cay}((\ZZ/N\ZZ)^*, S^\pm)$,
Proposition~\ref{prop:char_decomp} gives:
\begin{equation}\label{eq:link1}
  \mu_\chi = \sum_{s \in S^\pm} \chi(s).
\end{equation}
The generators $S$ are primes (specifically, the smallest primes
coprime to~$N$).  The Dirichlet L-function
$L(s,\chi) = \prod_p (1 - \chi(p) p^{-s})^{-1}$
is an Euler product over \emph{all} primes.  The eigenvalue
$\mu_\chi$ encodes precisely the Euler factors at the primes in~$S$.

As the primorial $N = p_1 \cdots p_k$ grows, the generators include
more primes, and $\mu_\chi$ samples more Euler factors of $L(s,\chi)$.
In the limit $k \to \infty$, the generating set exhausts the primes,
and the character eigenvalues determine the L-function completely.

This is verified: all $\varphi(N)$ eigenvalues at every ladder level
match the character sum formula to precision $< 5 \times 10^{-14}$,
and the principal character L-function reproduces $\zeta(s) \prod_{p
\mid N} (1 - p^{-s})$ to precision $< 5 \times 10^{-5}$.

\subsection{Link~2: The Graph Trace Formula Is the Explicit Formula
  \textnormal{[\textsc{verified}]}}

Theorem~\ref{thm:explicit} gives:
\[
  N_k = \Tr(B^k) = \sum_j \theta_j^k
  \qquad\text{and}\qquad
  \pi_\mathcal{G}(L) = \frac{1}{L} \sum_{d \mid L}
  \mu\!\left(\frac{L}{d}\right) N_d.
\]
This is the graph-theoretic analog of the explicit formula
\begin{equation}\label{eq:classical_explicit}
  \psi(x) = x - \sum_\rho \frac{x^\rho}{\rho} - \ln(2\pi)
  - \tfrac{1}{2}\ln(1 - x^{-2}),
\end{equation}
where the sum runs over nontrivial zeros $\rho$ of~$\zeta(s)$.  Both
formulas reconstruct prime-counting functions from spectral data.  The
difference: equation~\eqref{eq:classical_explicit} counts integer
primes from zeta zeros; the graph formula counts prime \emph{cycles}
from Hashimoto eigenvalues.

Verified to machine precision on all graphs
(Table~\ref{tab:pnt}).  On $\GL(2, \FF_3)$, the graph PNT ratio
$\pi_\mathcal{G}(L) / [(d-1)^L / L]$ converges to~1 by $L = 8$.

\subsection{Link~3: The Abelian Obstruction and Its Diagnosis
  \textnormal{[\textsc{verified}]}}
\label{sec:abelian_obstruction}

Links~1--2 show that the spectral machinery works, but on abelian
Cayley graphs it cannot go further.  The Davis Law
(Section~\ref{sec:davis_law}) diagnoses why:
\begin{enumerate}[nosep]
  \item All irreducible representations of $(\ZZ/N\ZZ)^*$ are
    one-dimensional.  A character can evaluate to $-1$ on every
    generator simultaneously.
  \item The CRT factorization~\eqref{eq:crt} allows characters to
    ``see'' only a depleted subset of factors.  When generators align
    on one factor (Section~\ref{sec:failure}), the character sum
    achieves maximum constructive interference: $\mu_\chi = -d$.
  \item This is the spectral analog of density depletion at an
    obstacle rim in superfluid flow.  The Ramanujan bound assumes
    uniform spectral density (the infinite tree), just as Landau's
    criterion assumes uniform fluid density.  Both fail for the same
    structural reason.
\end{enumerate}

The corrected bound $|\mu_\chi| \leq \rho_A(S,\chi)$ is exact
(Proposition~\ref{prop:corrected}), but it is a \emph{diagnostic},
not a remedy.  To reach Ramanujan and connect to the Riemann
Hypothesis, we need two things: a non-abelian group (to prevent
coherent alignment) and arithmetically chosen generators (to ensure
the Ramanujan bound holds as a theorem, not merely as a near-miss).

\subsection{Link~4: The LPS Construction Provides Both
  \textnormal{[\textsc{established}]}}

Lubotzky, Phillips, and Sarnak~\cite{lps1988} solve both problems
simultaneously.  For distinct primes $p, q$ with $p \equiv 1
\pmod{4}$ and $\bigl(\frac{p}{q}\bigr) = 1$:

\begin{enumerate}[nosep]
  \item \textbf{Non-abelian group:} The Cayley graph is built on
    $\PGL(2, \FF_q)$, whose irreducible representations have dimension
    up to $q-1$.  By the mechanism of
    Section~\ref{sec:gl2}, higher-dimensional representations
    suppress coherent alignment via destructive phase interference.
  \item \textbf{Arithmetic generators:} The $p+1$ generators are
    derived from integer quaternions
    \begin{equation}\label{eq:quaternion}
      \alpha = a_0 + a_1 i + a_2 j + a_3 k, \qquad
      a_0^2 + a_1^2 + a_2^2 + a_3^2 = p,
      \quad a_0 > 0 \text{ odd},
    \end{equation}
    mapped to $\PGL(2, \FF_q)$ via the standard embedding.  These are
    not ad hoc --- they are determined by the arithmetic of the
    quaternion algebra.
\end{enumerate}

\begin{theorem}[LPS~\cite{lps1988}, via Deligne~\cite{deligne1974}]
\label{thm:lps}
  The LPS Cayley graph satisfies $|\mu_k| \leq 2\sqrt{p}$ for all
  nontrivial eigenvalues.  This is a theorem, not a conjecture: the
  proof uses Deligne's resolution of the Ramanujan--Petersson
  conjecture for $\GL(2)$ over finite fields.
\end{theorem}

In the language of the Davis Law: the LPS construction
\emph{eliminates the hard obstacle}.  The quaternion generators are
arithmetically balanced (no QR alignment pathology), and the
non-abelian representation structure prevents density depletion.
The $K_{\mathrm{eff}}$ diagnostic of Section~\ref{sec:davis_law}
predicts this: soft obstacle, low $K_{\mathrm{eff}}$, Ramanujan
achieved.

\subsection{Link~5: The LPS Adjacency Operator Is a Hecke Operator
  \textnormal{[\textsc{established}]}}

This is the critical identification.  The adjacency operator of the
LPS graph is not merely \emph{analogous to} a Hecke operator --- it
\emph{is} one:
\begin{equation}\label{eq:lps_hecke}
  A_{\mathrm{LPS}} = T_p
\end{equation}
where $T_p$ is the Hecke operator at the prime~$p$, acting on
$L^2(\PGL(2, \FF_q))$.  Its eigenvalues are the Hecke eigenvalues
$a_p(\pi)$ for each automorphic representation~$\pi$ of
$\PGL(2)$.

This means:
\begin{itemize}[nosep]
  \item Each eigenvalue of the LPS graph is labeled by an automorphic
    representation~$\pi$.
  \item The Ramanujan bound $|a_p(\pi)| \leq 2\sqrt{p}$ is the
    Ramanujan--Petersson conjecture for that representation.
  \item The ``prime cycles'' of the LPS graph correspond to conjugacy
    classes in $\PGL(2, \FF_q)$ determined by the quaternion algebra's
    prime ideal structure.
\end{itemize}

\emph{Cuspidal extraction on LPS graphs.}  Merging the LPS construction
(Link~4) with the isotypic decomposition machinery of
Section~\ref{sec:isotypic}, we decompose the Hecke operator $T_p$ on
$\PSL(2, \FF_q)$ into representation sectors using class-sum operators
from the center of the group algebra.  For $p = 5$, $q = 11$: three
cuspidal Hecke eigenvalues are identified at $a_\pi \in \{-3.00,
\pm 1.73\}$, with maximum $|a_\pi| / 2\sqrt{5} = 0.671$; for
$p = 17$, $q = 13$: maximum cuspidal ratio $0.723$.  In every case the
cuspidal eigenvalues satisfy the Ramanujan bound with substantial
margin, consistent with the fact that these are exactly the eigenvalues
governed by Deligne's theorem.  Each cuspidal $a_p(\pi)$ feeds into
a local Euler factor
$L_p(s, \pi) = (1 - a_p(\pi)\, p^{-s} + p^{1-2s})^{-1}$,
demonstrating the direct path from finite graph computation to
L-function data.

\subsection{Link~6: Jacquet--Langlands Lifts to $\GL(2)$
  \textnormal{[\textsc{established}]}}

The Jacquet--Langlands correspondence~\cite{jacquet_langlands1970}
provides the bridge from the quaternion algebra (where LPS lives) to
the standard setting of $\GL(2)$ (where L-functions live):

\begin{equation}\label{eq:jl}
  \pi' \text{ on } D^\times / \FF_q^\times
  \;\;\xleftrightarrow{\;\;\text{JL}\;\;}\;\;
  \pi \text{ on } \GL(2, \FF_q)
\end{equation}

where $D$ is the quaternion algebra and $\pi, \pi'$ are automorphic
representations.  The correspondence preserves L-functions:
$L(s, \pi') = L(s, \pi)$.  In particular, their Hecke eigenvalues
match: $a_p(\pi') = a_p(\pi)$.

For the trivial central character, the L-functions of $\GL(2)$
automorphic representations include, via the Langlands program,
L-functions attached to modular forms.  The Riemann zeta function
$\zeta(s)$ itself appears as the L-function of the trivial
representation of $\GL(1)$, which embeds into $\GL(2)$ via the
Eisenstein series.  The \emph{cuspidal} representations --- those
that do not arise from $\GL(1)$ --- carry the genuinely new spectral
information.

\subsection{Link~7: The Adelic Limit
  \textnormal{[\textsc{open}]}}

The chain so far works at a single prime~$q$.  To recover the full
Riemann zeta function, we need the adelic picture.  The Satake
isomorphism identifies:
\begin{equation}\label{eq:satake}
  a_p(\pi) = \alpha_p + \alpha_p^{-1}
\end{equation}
where $\alpha_p$ is the Satake parameter of $\pi$ at~$p$.  The
automorphic L-function is
\begin{equation}\label{eq:automorphic_L}
  L(s, \pi) = \prod_p \frac{1}{(1 - \alpha_p\, p^{-s})
  (1 - \alpha_p^{-1}\, p^{-s})}.
\end{equation}
The Ramanujan--Petersson conjecture states $|\alpha_p| = 1$ for
cuspidal~$\pi$, which forces the zeros of $L(s,\pi)$ to lie on
$\operatorname{Re}(s) = 1/2$.

The remaining program is:
\begin{enumerate}[nosep]
  \item \textbf{Build LPS graphs} for increasing $q$ and verify
    Theorem~\ref{thm:lps} computationally.
  \item \textbf{Label eigenvalues} by representation type: principal
    series ($\dim = q-1$ or $q+1$), Steinberg ($\dim = q$), cuspidal
    ($\dim = q-1$).  \emph{This step is now computationally realized
    for $\GL(2, \FF_p)$ with ad hoc generators at $p \leq 5$
    (Section~\ref{sec:isotypic}).  Extending to LPS quaternion
    generators is the next milestone.}
  \item \textbf{Extract Hecke eigenvalues} $a_p(\pi)$ from cuspidal
    representations and verify they satisfy the Weil bound
    $|a_p| \leq 2\sqrt{p}$.  \emph{Verified at $p = 3$ with all
    cuspidal eigenvalues Ramanujan
    (Table~\ref{tab:cuspidal_hecke}).}
  \item \textbf{Vary $p$ and $q$:} for a fixed cuspidal~$\pi$,
    compute $a_p(\pi)$ at multiple primes~$p$ and reconstruct the
    partial Euler product~\eqref{eq:automorphic_L}.
  \item \textbf{Take $q \to \infty$:} show that the graph prime cycle
    distribution $\pi_\mathcal{G}(L)$ converges to the integer prime
    distribution $\pi(x)$ under the substitution
    $x = (d-1)^{L/2}$.
\end{enumerate}

\subsection{Link~8: The Archimedean Completion
  \textnormal{[\textsc{verified}]}}
\label{sec:archimedean}

The partial Euler product $\prod_{p \leq X} L_p(s, \pi)$ converges
absolutely only for $\operatorname{Re}(s) > 1$.  To reach the critical
line, one needs the \emph{completed} L-function
\begin{equation}\label{eq:completed}
  \Lambda(s, f) = N^{s/2} \, \Gamma_\RR(s + \mu_1) \,
  \Gamma_\RR(s + \mu_2) \, L(s, f),
\end{equation}
which satisfies a functional equation
$\Lambda(s) = \varepsilon \, \Lambda(k - s)$.
The parameters $N$, $k$, $\mu_1$, $\mu_2$, $\varepsilon$ constitute
the \emph{archimedean data} --- the contribution of the infinite
place of~$\QQ$.

For the LPS construction on $\PSL(2, \FF_{11})$, the quaternion algebra
$B$ ramified at $\{11, \infty\}$ determines all archimedean parameters
with no free choices:

\begin{enumerate}[nosep]
  \item \textbf{Weight.}  $B \otimes_\QQ \RR \cong \mathbb{H}$
    (Hamilton quaternions), so $B$ is \emph{definite}.  Under the
    Jacquet--Langlands correspondence, automorphic forms on~$B^\times$
    map to \emph{holomorphic} modular forms.  The trivial $K$-type on
    $B^\times(\RR) = \mathbb{H}^\times$ maps to the discrete series
    $D_2$, forcing $k = 2$.
  \item \textbf{Conductor.}  The discriminant of~$B$ is~$11$.  At
    $p = 11$ the local component is Steinberg (conductor exponent~$1$);
    at $p \neq 11$ it is unramified (exponent~$0$).  Therefore $N = 11$.
  \item \textbf{Gamma factors.}  For a weight-$2$ holomorphic newform:
    $L_\infty(s) = \Gamma_\CC(s) = (2\pi)^{-s}\, \Gamma(s) =
    \Gamma_\RR(s) \cdot \Gamma_\RR(s+1)$, giving $\mu_1 = 0$,
    $\mu_2 = 1$.
  \item \textbf{Root number.}  $\varepsilon = (-1)^{k/2} \cdot w_{11}
    = (-1) \cdot (-1) = +1$, where $w_{11} = -1$ is the Atkin--Lehner
    eigenvalue at the prime of split multiplicative reduction.
  \item \textbf{Uniqueness.}  $\dim S_2(\Gamma_0(11))^{\mathrm{new}}
    = \mathrm{genus}(X_0(11)) = 1$.  There is exactly one newform at
    this level and weight: $f = \eta(\tau)^2 \eta(11\tau)^2$, the
    modular form of the elliptic curve $y^2 + y = x^3 - x^2 - 10x - 20$
    (Cremona label 11a1, LMFDB label 11.2.a.a).
\end{enumerate}

\begin{table}[ht]
\centering
\renewcommand{\arraystretch}{1.1}
\begin{tabular}{lll}
  \toprule
  Parameter & Value & Determined by \\
  \midrule
  Conductor $N$ & $11$ & $\mathrm{disc}(B) = 11$ \\
  Weight $k$ & $2$ & $B$ definite $\to$ discrete series $D_2$ \\
  $\mu_1, \mu_2$ & $0, 1$ & $\Gamma_\CC(s) = \Gamma_\RR(s) \cdot \Gamma_\RR(s+1)$ \\
  Root number $\varepsilon$ & $+1$ & $(-1)^{k/2} \cdot w_{11} = (-1)(-1)$ \\
  Analytic rank & $0$ & $L(1, f) \approx 0.2538 \neq 0$ \\
  $\dim S_2(\Gamma_0(11))^{\mathrm{new}}$ & $1$ & $\mathrm{genus}(X_0(11)) = 1$ \\
  \bottomrule
\end{tabular}
\caption{Archimedean parameters of the completed L-function, all forced
  by the quaternion algebra $B$ ramified at $\{11, \infty\}$.  No
  parameter is fitted; the algebra determines the infinite place.}
\label{tab:archimedean}
\end{table}

We verified the functional equation numerically.  With $2000$ Fourier
coefficients of~$f$ computed from the $\eta$-product formula, the ratio
$\Lambda(s)/\Lambda(2-s) \to 1$ as $s \to 1^+$ from above
(Table~\ref{tab:functional_eqn}).  The Dirichlet series does not
converge for $\operatorname{Re}(s) < 1$; this is precisely why the
archimedean completion is required.  It provides the bridge from
convergent partial Euler products ($\operatorname{Re}(s) > 1$) to the
critical line ($\operatorname{Re}(s) = 1/2$ in analytic normalization).

\begin{table}[ht]
\centering
\renewcommand{\arraystretch}{1.1}
\begin{tabular}{ccrrr}
  \toprule
  $s$ & $2-s$ & $\Lambda(s)$ & $\Lambda(2-s)$ & $\Lambda(s)/\Lambda(2-s)$ \\
  \midrule
  $1.01$ & $0.99$ & $0.11234$ & $0.10816$ & $1.039$ \\
  $1.02$ & $0.98$ & $0.11417$ & $0.10577$ & $1.079$ \\
  $1.05$ & $0.95$ & $0.11880$ & $0.09718$ & $1.222$ \\
  $1.10$ & $0.90$ & $0.12431$ & $0.07646$ & $1.626$ \\
  \bottomrule
\end{tabular}
\caption{Functional equation $\Lambda(s) = \varepsilon \Lambda(2-s)$
  with $\varepsilon = +1$, tested with $2000$-term Dirichlet series.
  The ratio approaches~$1$ as $s \to 1^+$; deviation at larger~$s$
  reflects Dirichlet non-convergence at $2 - s < 1$, confirming that
  the archimedean completion is necessary.}
\label{tab:functional_eqn}
\end{table}

The Sato--Tate distribution of the known newform~$f$ matches the
semicircle law: across $302$ primes $p \leq 2000$, the mean
Sato--Tate angle is $\bar\theta = 1.565$ (predicted: $\pi/2 = 1.571$).
Restricting to the $147$ LPS-eligible primes ($p \equiv 1 \pmod{4}$,
$p \leq 2000$) gives $\bar\theta = 1.561$, confirming equidistribution
by the Chebotarev density theorem.


\subsection{Summary: The Complete Chain}

\begin{equation}\label{eq:chain}
  \boxed{
  \begin{aligned}
    &\text{Graph eigenvalue } \mu_k \\
    &\quad\xrightarrow{\text{Link 5: LPS = Hecke}}
    \text{Hecke eigenvalue } a_p(\pi) \\
    &\quad\xrightarrow{\text{Link 6: Jacquet--Langlands}}
    \text{Automorphic } L(s, \pi) \text{ (partial Euler product)} \\
    &\quad\xrightarrow{\text{Link 8: Archimedean completion}}
    \Lambda(s, f) = N^{s/2}\, \Gamma_\RR(s)\, \Gamma_\RR(s{+}1) \cdot L(s, f)
  \end{aligned}
  }
\end{equation}

The completed L-function $\Lambda(s, f)$ satisfies
$\Lambda(s) = \varepsilon\, \Lambda(2 - s)$ with $\varepsilon = +1$,
$N = 11$, $k = 2$.  This functional equation provides the analytic
continuation to all of~$\CC$ and forces the nontrivial zeros onto the
critical line $\operatorname{Re}(s) = 1$ (arithmetic normalization)
$= \operatorname{Re}(s) = 1/2$ (analytic normalization).

For the unique cuspidal representation at level~$11$, the
Ramanujan--Petersson conjecture (proven over finite fields by Deligne,
open over~$\QQ$) together with the completed functional equation
constrains the zero distribution completely.  The Hilbert--P\'olya
operator in this framework is the Hecke operator $T_p$ on the cuspidal
sector: its self-adjointness is guaranteed by the Petersson inner
product, and its eigenvalues $a_p(\pi)$ determine the zeros of
$\Lambda(s, f)$ via the Euler product.

\subsection{What the Davis Law Contributes to This Program}

The chain~\eqref{eq:chain} is known, at least in outline, to the
automorphic forms community.  What is new in this work is the
\emph{diagnostic layer} provided by the Davis Law.

\begin{enumerate}
  \item \textbf{Failure prediction.}  The Davis Law predicts
    \emph{which} Cayley graphs will fail the Ramanujan bound and
    \emph{why}, before computing a single eigenvalue.  The CRT
    coverage analysis (Table~\ref{tab:coverage}) and QR alignment
    check are computable from generators alone.  No existing result in
    the Ramanujan graph literature provides this level of a priori
    diagnostic.

  \item \textbf{Generator design.}  The corrected bound $|\mu_\chi|
    \leq \rho_A(S,\chi)$ tells the graph designer exactly what to
    optimize: minimize $K_{\mathrm{eff}}$ by balancing quadratic
    residue classes across all CRT factors.  The LPS quaternion
    construction achieves this automatically via the arithmetic of
    $\sum a_i^2 = p$ representations, but the Davis Law explains
    \emph{why} it works and suggests that other constructions
    satisfying the same coverage conditions would also succeed.

  \item \textbf{The superfluid correspondence.}  The parallel between
    spectral density depletion and superfluid density depletion
    (Table~\ref{tab:parallel}) is not merely metaphorical.  Both
    systems exhibit:
    \begin{itemize}[nosep]
      \item a uniform-background bound that fails at finite obstacles
        (Landau / Ramanujan),
      \item a local diagnostic ($c_s(r^*)/A(r^*)$ / $\rho_A(S,\chi)$)
        that predicts the corrected threshold exactly,
      \item a classification of obstacles as ``hard'' (impenetrable,
        high $K_{\mathrm{eff}}$: abelian CRT / rigid body) or ``soft''
        (penetrable, low $K_{\mathrm{eff}}$: non-abelian reps /
        porous medium).
    \end{itemize}
    Whether this correspondence extends deeper --- whether there is a
    precise mathematical functor between the category of arithmetic
    Cayley graphs and the category of obstacle geometries in quantum
    fluids --- is an open question that the present framework is
    designed to investigate.  The fact that the corrected bounds are
    \emph{exact} in both settings suggests the parallel is structural,
    not accidental.

  \item \textbf{The geodesic interpretation.}  In the superfluid
    setting, the critical velocity $v_c$ determines which flow lines
    (geodesics of the fluid) can propagate past the obstacle without
    nucleating excitations.  In the spectral setting, the Ramanujan
    bound determines which ``geodesics'' (prime cycles in the Cayley
    graph) can be reconstructed from the spectrum without spectral
    aliasing.  The Davis Law identifies the local geometry that governs
    this transition.  Pursuing this interpretation on LPS graphs ---
    where prime cycles correspond to actual arithmetic data via
    Links~5--6 --- is the most direct path from the present results to
    the Hilbert--P\'olya conjecture.
\end{enumerate}


%% ====================================================================
\section{Discussion and Assessment}\label{sec:discussion}
%% ====================================================================

\subsection{What the Primorial Ladder Proves}

\begin{figure}[ht]
\centering
\includegraphics[width=0.85\textwidth]{fig7_riemann_comparison.png}
\caption{Rescaled manifold eigenvalues (colored) vs.\ Riemann zeta
  zeros $t_n$ (black stars).  As the primorial level increases, the
  eigenvalue staircase more closely approximates the distribution of
  Riemann zeros, motivating the Hilbert--P\'olya program.  The
  convergence is suggestive but not sufficient --- connecting to actual
  zeros requires the automorphic machinery of Section~\ref{sec:lps}.}
\label{fig:riemann_comparison}
\end{figure}

The ladder is not an experiment that succeeded or failed.  It is a
logical argument in four steps:
\begin{enumerate}[nosep]
  \item \textbf{Levels~0--2:} The spectral machinery works.  Character
    decomposition, Ihara zeta, heat-zeta bridge, Fisher geometry,
    thermodynamics, Hodge theory, RG flow --- seven cross-branch
    identities verified to machine precision.
  \item \textbf{Level~3:} Abelian Cayley graphs with CRT structure
    cannot be Ramanujan at scale.  The mechanism is diagnosed exactly:
    quadratic residue alignment on depleted CRT factors
    (Section~\ref{sec:failure}).
  \item \textbf{$\GL(2)$:} Non-abelian representation theory resolves
    the obstruction.  GUE statistics confirm the correct universality
    class (Section~\ref{sec:gl2}).
  \item \textbf{Trace formula:} The spectral--arithmetic bridge is
    verified: eigenvalues reconstruct prime cycle counts via the
    Hashimoto operator (Section~\ref{sec:trace}).
\end{enumerate}
Each step forces the next.  The failure at Level~3 is not a setback
--- it is the datum that reveals the mechanism and points to the
resolution.

\subsection{What Remains}

The seven links of Section~\ref{sec:lps} are in three states:

\begin{table}[ht]
\centering
\renewcommand{\arraystretch}{1.2}
\begin{tabular}{clll}
  \toprule
  Link & Content & Status & Evidence \\
  \midrule
  1 & Eigenvalues $=$ partial L-functions & \textsc{verified}
    & All levels, error $< 10^{-13}$ \\
  2 & Graph trace formula & \textsc{verified}
    & All graphs, error $< 10^{-7}$ \\
  3 & Abelian obstruction diagnosed & \textsc{verified}
    & Davis Law, ratio $= 1.000$ \\
  4 & LPS construction & \textsc{established}
    & \cite{lps1988}, not yet implemented \\
  5 & LPS $=$ Hecke operator & \textsc{established}
    & \cite{lps1988, lubotzky1994} \\
  6 & Jacquet--Langlands & \textsc{established}
    & \cite{jacquet_langlands1970} \\
  7a & Rep.\ labeling ($p \leq 5$) & \textsc{verified}
    & Table~\ref{tab:isotypic}, error $< 10^{-14}$ \\
  7b & Cuspidal Hecke eigenvalues & \textsc{verified}
    & Table~\ref{tab:cuspidal_hecke}, $p=3$ Ramanujan \\
  7c & Adelic limit $q \to \infty$ & \textsc{open}
    & Requires passage to adeles \\
  8 & Archimedean completion & \textsc{verified}
    & Table~\ref{tab:archimedean}: $N{=}11$, $k{=}2$, $\varepsilon{=}{+}1$ \\
  9 & Functional equation & \textsc{verified}
    & Table~\ref{tab:functional_eqn}, $\Lambda(s) = \Lambda(2{-}s)$ \\
  \bottomrule
\end{tabular}
\caption{Status of each link in the chain from graph primes to
  integer primes.  Links~7a--7b advance the former ``open'' Link~7
  by computationally verifying representation labeling and Hecke
  eigenvalue extraction for $\GL(2, \FF_p)$ with $p \leq 5$.
  Links~8--9 close the archimedean place: the quaternion algebra
  determines $N$, $k$, $\mu$, $\varepsilon$ with no free parameters.}
\label{tab:links}
\end{table}

The first three links are original computational results of this work.
Links~4--6 are theorems in the literature that have not yet been
implemented in our pipeline.  Links~7a--7b are new computational
results: the isotypic decomposition of $\GL(2, \FF_p)$ Cayley graphs
labels \emph{every} eigenvalue by representation type
(Table~\ref{tab:isotypic}) and extracts the cuspidal Hecke eigenvalues
that constitute the Hilbert--P\'olya sector
(Table~\ref{tab:cuspidal_hecke}).  Links~8--9 complete the L-function
at the infinite place: the quaternion algebra $B$ ramified at
$\{11, \infty\}$ forces weight $k = 2$, conductor $N = 11$, and
root number $\varepsilon = +1$, uniquely identifying the newform
$f = \eta(\tau)^2\eta(11\tau)^2$ of the elliptic curve 11a1
(Table~\ref{tab:archimedean}).

Link~7c --- the adelic limit ---
is the genuine open problem:
it requires the passage from finite $\FF_q$ to the adeles, which is
the deepest part of the Langlands program.

\subsection{The Role of Novel Diagnostics}

Classical approaches to the Hilbert--P\'olya conjecture
(Berry--Keating~\cite{berry_keating1999},
Connes~\cite{connes1999}) work top-down: start with $\zeta(s)$
and search for the operator.  The present approach works bottom-up:
start with finite computable models, verify the structural
relationships, diagnose failures, and let the failures point toward
the resolution.

The Davis Law of Ramanujan is a product of this bottom-up approach.
It could not have been discovered by studying $\zeta(s)$ directly,
because it concerns the \emph{local geometry of the generating set}
--- a feature invisible in the limit but decisive at finite scale.
The superfluid correspondence emerged from the same diagnostic
methodology, and it remains to be seen whether it can be elevated
from a structural parallel to a precise mathematical equivalence.

The framework is designed to accommodate such discoveries.  The
five-branch structure (spectral, thermodynamic, information-geometric,
cohomological, renormalization) provides multiple independent
diagnostic channels.  Each branch verified the same arithmetic
structure from a different angle, and each produced its own
constraints on the path forward.  Future phases of this research will
determine whether these constraints collectively force the existence
of the Hilbert--P\'olya operator, or reveal a deeper obstruction that
reshapes the question.

\subsection{Relation to the Geometry of Sameness}

The structure of the seven-link chain can be understood through the
lens of the \emph{Geometry of Sameness} framework~\cite{davis2025sameness},
which formalizes the duality between ``translator-based'' and
``manifold-based'' views of a fixed semantic sameness structure~$S$.

In the present setting, the hidden variable $S$ is the distribution
of primes (equivalently, the zeta zeros).  The \emph{translator view}
($\SamTrans$) is the finite graph: computable eigenvalues, local
generators, spectral statistics.  The \emph{manifold view}
($\SamGeom$) is the L-function: analytic continuation, functional
equation, zeros on the critical strip.  The chain of
Section~\ref{sec:lps} is the functor $F_S : \SamTrans \to \SamGeom$
that maps graph eigenvalues to L-function data; the LPS construction
is the inverse functor $G_S$ that derives graph generators from the
arithmetic structure of quaternion algebras.  The archimedean
completion (Link~8) is the point where the functor crosses from finite
to infinite: the algebra provides the gamma factors that complete the
analytic continuation.

The Davis Law plays the role of the \emph{Error Budget Transfer
Theorem}: it quantifies the slack between the finite (translator)
computation and the infinite (manifold) structure.  The ``smooth
chartability'' condition --- which gates whether the functor $F_S$
succeeds --- corresponds to the Ramanujan property: precisely when the
generating set is arithmetically balanced (LPS quaternion generators,
soft obstacle), the eigenvalue-to-L-function map is well-defined with
controlled error.  When chartability fails (abelian CRT, hard
obstacle), the functor breaks down and the Davis Law diagnoses the
failure mode exactly.

This dual-realization perspective explains why the bottom-up
approach works: the ``translator'' side is where computation happens,
while the ``manifold'' side is where the theorems live.  The functors
$F_S$ and $G_S$ ensure that verified properties of the computable
objects transfer to the analytic objects with bounded slack, and vice
versa.


%% ====================================================================
\section{Conclusion}\label{sec:conclusion}
%% ====================================================================

The primorial ladder establishes a multi-branch computational
framework for investigating the spectral geometry of arithmetic groups
in the context of the Hilbert--P\'olya conjecture.

The principal new results are:

\textbf{The Davis Law of Ramanujan}
(Proposition~\ref{prop:corrected}): the classical bound $|\mu| \leq
2\sqrt{d-1}$ is the ``Landau criterion'' of spectral graph theory,
valid only under the assumption of uniform spectral density.  When
this uniformity breaks --- through CRT-factored constructive
interference in the abelian setting, or through density depletion at
an obstacle rim in the superfluid setting --- the bound fails, and the
corrected bound is character-dependent and exact.

\textbf{The seven-link chain} (Section~\ref{sec:lps}): a precise
logical path from graph eigenvalues to automorphic L-function zeros,
with three links computationally verified, three established in the
literature, and one open.  The chain makes explicit how and why the
prime cycles of arithmetic Cayley graphs can, in principle, recover
the distribution of integer primes.

\textbf{The archimedean completion} (Section~\ref{sec:archimedean}):
the quaternion algebra $B$ ramified at $\{11, \infty\}$ determines the
infinite place of~$\QQ$ with no free parameters --- conductor $N = 11$,
weight $k = 2$, gamma factors $\Gamma_\RR(s) \cdot \Gamma_\RR(s+1)$,
root number $\varepsilon = +1$ --- uniquely identifying the completed
L-function as that of the elliptic curve 11a1.  The partial Euler
product (finite places from LPS data) and the archimedean factor
(forced by the algebra) together yield a complete L-function satisfying
a functional equation, closing the gap from convergent data to the
critical line.

\textbf{The diagnostic methodology}: the five-branch framework
(spectral, thermodynamic, information-geometric, cohomological,
renormalization) provides independent verification channels and
produces constraints --- like the Davis Law --- that are invisible
from any single branch alone.  The superfluid correspondence emerged
from this methodology and remains the most suggestive open direction.

\medskip

\emph{This is not a proof of the Riemann Hypothesis.  It is a
research program with verified computational foundations, a precise
failure mechanism, a new diagnostic law, and a complete chain connecting
finite spectral data to the zeros of automorphic L-functions.  Of the
eleven links in that chain (Table~\ref{tab:links}), seven are
computationally verified, three are theorems in the literature, and
one --- the adelic limit --- is open.  The L-function is complete at
all places of~$\QQ$; what remains is the passage from a single
finite field $\FF_q$ to the adeles.}



%% ====================================================================
\section*{Notation}
%% ====================================================================

\begin{tabular}{ll}
  $\zeta(s)$ & Riemann zeta function \\
  $\zeta_\mathcal{G}(u)$ & Ihara zeta function of graph $\mathcal{G}$ \\
  $L(s,\chi)$ & Dirichlet L-function for character~$\chi$ \\
  $\rho$ & Non-trivial zero of $\zeta(s)$ \\
  $\mu_j$ & Adjacency eigenvalue ($= d - \lambda_j$) \\
  $\lambda_j$ & Laplacian eigenvalue \\
  $\theta_j$ & Hashimoto (non-backtracking) eigenvalue \\
  $B$ & Hashimoto non-backtracking matrix \\
  $\Theta(t) = \Tr(e^{-tL})$ & Heat trace \\
  $Z(\beta) = \Tr(e^{-\beta L})$ & Partition function \\
  $g_F$ & Fisher information metric \\
  $\Delta^F$ & Fisher Laplacian \\
  $\kappa$ & Ollivier--Ricci curvature \\
  $\xi$ & Correlation length ($= 1/\lambda_1$) \\
  $C_{\mathrm{sp}}$ & Spectral capacity ($= \lambda_1 D^2$) \\
  GUE & Gaussian Unitary Ensemble \\
  GRH & Generalized Riemann Hypothesis \\
  CRT & Chinese Remainder Theorem \\
  QR/QNR & Quadratic residue / non-residue \\
  LPS & Lubotzky--Phillips--Sarnak \\
\end{tabular}


%% ====================================================================
%% REFERENCES
%% ====================================================================

\begin{thebibliography}{99}

\bibitem{davis2024geometry}
B.~R.~Davis,
\emph{The Geometry of Sameness},
Amazon KDP, December 2024.

\bibitem{davis2026landau}
B.~R.~Davis,
``Compressible sonic onset and the Davis Law $C = \tau/K$
for Bose--Einstein condensate flow past finite obstacles,''
preprint, February 2026.

\bibitem{lps1988}
A.~Lubotzky, R.~Phillips, and P.~Sarnak,
``Ramanujan graphs,''
\emph{Combinatorica} \textbf{8} (1988), no.~3, 261--277.

\bibitem{margulis1988}
G.~A.~Margulis,
``Explicit group-theoretical constructions of combinatorial schemes
and their application to the design of expanders and concentrators,''
\emph{Problemy Peredachi Informatsii} \textbf{24} (1988), no.~1,
51--60.

\bibitem{deligne1974}
P.~Deligne,
``La conjecture de Weil.~I,''
\emph{Inst.\ Hautes \'Etudes Sci.\ Publ.\ Math.} \textbf{43}
(1974), 273--307.

\bibitem{selberg1956}
A.~Selberg,
``Harmonic analysis and discontinuous groups in weakly symmetric
Riemannian spaces with applications to Dirichlet series,''
\emph{J.~Indian Math.~Soc.} \textbf{20} (1956), 47--87.

\bibitem{jacquet_langlands1970}
H.~Jacquet and R.~P.~Langlands,
\emph{Automorphic forms on $\mathrm{GL}(2)$},
Lecture Notes in Mathematics, vol.~114, Springer-Verlag, 1970.

\bibitem{ihara1966}
Y.~Ihara,
``On discrete subgroups of the two by two projective linear group
over $\mathfrak{p}$-adic fields,''
\emph{J.~Math.~Soc.~Japan} \textbf{18} (1966), 219--235.

\bibitem{bass1992}
H.~Bass,
``The Ihara--Selberg zeta function of a tree lattice,''
\emph{Internat.~J.~Math.} \textbf{3} (1992), no.~6, 717--797.

\bibitem{hashimoto1989}
K.~Hashimoto,
``Zeta functions of finite graphs and representations of $p$-adic
groups,''
in \emph{Automorphic Forms and Geometry of Arithmetic Varieties},
Adv.~Stud.~Pure Math., vol.~15, Academic Press, 1989, pp.~211--280.

\bibitem{montgomery1973}
H.~L.~Montgomery,
``The pair correlation of zeros of the zeta function,''
\emph{Proc.~Sympos.~Pure Math.} \textbf{24} (1973), 181--193.

\bibitem{odlyzko1987}
A.~M.~Odlyzko,
``On the distribution of spacings between zeros of the zeta
function,''
\emph{Math.~Comp.} \textbf{48} (1987), no.~177, 273--308.

\bibitem{berry_keating1999}
M.~V.~Berry and J.~P.~Keating,
``The Riemann zeros and eigenvalue asymptotics,''
\emph{SIAM Rev.} \textbf{41} (1999), no.~2, 236--266.

\bibitem{connes1999}
A.~Connes,
``Trace formula in noncommutative geometry and the zeros of the
Riemann zeta function,''
\emph{Selecta Math.~(N.S.)} \textbf{5} (1999), no.~1, 29--106.

\bibitem{lubotzky1994}
A.~Lubotzky,
\emph{Discrete Groups, Expanding Graphs and Invariant Measures},
Progress in Mathematics, vol.~125, Birkh\"auser, 1994.

\bibitem{stark_terras2000}
H.~M.~Stark and A.~A.~Terras,
``Zeta functions of finite graphs and coverings, Part~II,''
\emph{Adv.\ Math.} \textbf{154} (2000), no.~1, 132--195.

\bibitem{davis2025sameness}
B.~R.~Davis,
``The Geometry of Sameness: An $\varepsilon$-equivalence of translation
and distance,''
Manuscript, 2025.

\end{thebibliography}

\end{document}
